\chapter{Ethics, Law and the Person of Interest Question}
\label{ch:ethics}

The project was inspired by a fictional system of mass surveillance. That
inspiration is a reason for more care, not less. This chapter states what
the system will not do and why. It summarises the legal framework that
applies in India and the European Union, sets out the design rules that
follow, and drafts the ethics statement the paper will carry.

\section{What the Machine does that this system will not}

The Machine of \emph{Person of Interest} combines three capabilities beyond
detection~\citep{poi2011}.
\begin{itemize}
  \item It \emph{identifies} individuals and follows them across cameras.
  \item It \emph{profiles} them from their movements and associations.
  \item It \emph{predicts} who will commit or suffer a crime.
\end{itemize}
Each of these is, in law, a distinct and heavily regulated activity: remote
biometric identification, profiling and predictive risk assessment.

The system in this dossier is a perception layer. It detects and tracks
objects, and it tracks people only as anonymous moving boxes. It does not
recognise faces, it builds no identities and it makes no predictions about
individuals. The gap between the two is not a technical limitation waiting
to be closed. It is a boundary the project keeps on purpose.

\section{The legal framework}

\paragraph{European Union}
The AI Act, Regulation (EU) 2024/1689, has prohibited a set of practices since
2~February~2025~\citep{euaiact2024}. Article~5 bans four that matter here.
\begin{itemize}
  \item Assessing a person's risk of committing a criminal offence based
        solely on profiling or personality traits.
  \item Building facial-recognition databases by untargeted scraping of
        images from the internet or CCTV.
  \item Inferring sensitive characteristics, such as political opinion or
        religion, from biometric data.
  \item Real-time remote biometric identification in publicly accessible
        spaces for law enforcement, apart from three narrow exceptions that
        each require prior authorisation.
\end{itemize}
Other remote biometric identification is classed as high-risk. The
application date of those obligations was postponed to December 2027 by
later amendment\unv. Systems developed solely for scientific research are
outside the Act's scope, but their deployment is not.

The General Data Protection Regulation adds two further constraints. It
treats biometric data used for unique identification as a special category.
It requires an impact assessment before systematic monitoring of publicly
accessible areas on a large scale~\citep{gdpr2016}.

\paragraph{India}
The Digital Personal Data Protection Act, 2023, and its Rules, notified in
November 2025, come into force in phases. The core obligations apply from
May 2027\unv~\citep{dpdp2023,dpdprules2025}. Three provisions matter here.
\begin{itemize}
  \item Processing personal data, including images of identifiable people,
        requires a lawful basis, normally consent.
  \item The exemption for research and statistical purposes applies only
        when no decision is taken about a specific individual and prescribed
        safeguards are met.
  \item Tracking or behavioural monitoring of children is prohibited.
\end{itemize}
Beneath the Act sits the Supreme Court's recognition of privacy as a
fundamental right in \emph{Puttaswamy}. Under that judgment, any
surveillance must be lawful, necessary and proportionate~\citep{puttaswamy2017}.

\section{Design rules}
\label{sec:design-rules}

The rules below apply in every phase. They are written so that a reviewer,
an ethics committee or a future user can check them.

\begin{enumerate}
  \item No face recognition, re-identification of people, watch-lists or
        identity databases are built or used.
  \item Faces in any footage the project records are blurred at ingest,
        before storage.
  \item Processing happens on the device where possible. Only detections
        and aggregates leave it.
  \item Phase~4 outputs are aggregates and rule-based events about places,
        not profiles of people.
  \item Raw footage the project records is deleted on a fixed schedule.
        Retention periods are written down in advance.
  \item Recording for the project happens only with consent or clear
        signage, and never in places where children are likely to be
        present.
  \item Approval from the university's ethics committee is obtained before
        any recording of people.
  \item The withdrawn datasets of Table~\ref{tab:withdrawn} are never used.
  \item Released models carry a model card that states intended uses and
        excluded uses. Identification and predictive policing are among the
        excluded uses.
  \item Access to live feeds and logs is restricted and itself logged.
\end{enumerate}

\section{Uses worth building for}

The same perception layer serves purposes that need no identification at
all, and the paper should present these as its motivation.
\begin{itemize}
  \item Search and rescue at night or at sea, where thermal imaging finds
        people who cannot be seen.
  \item Detection of small drones near airports and other sensitive sites.
        Drones are one of RGBT-Tiny's seven classes.
  \item Wildlife monitoring and the detection of poaching.
  \item Traffic safety in fog and darkness.
  \item Perimeter safety on industrial sites.
  \item Disaster response from drones.
\end{itemize}

\section{Draft ethics and dual-use statement}

The paper will include a statement along the following lines.

\begin{quote}
\itshape
This work develops a detector for small objects in visible and thermal video.
It does not perform face recognition, re-identification or any form of
biometric identification, and it was not evaluated for such use. All
datasets were used under their research licences. No dataset withdrawn for
privacy reasons was used. Footage recorded by the authors was collected with
the approval of the institutional ethics committee, with faces blurred at
capture. We recognise that detection and tracking can be components of
surveillance systems. We therefore release the model with a card that
excludes identification and predictive policing, and we encourage
deployments that comply with the EU AI Act and India's Digital Personal Data
Protection Act.
\end{quote}
